\section{Candidate actions and doubles constraints}
\label{sec:actions}

\subsection{Actions are complete player decisions}
For an actionable request $r$, let $\A(r)$ be the list returned by
\code{legal_choices(r)}. In doubles, one candidate is a complete pair
$a=(u_1,u_2)$. Each $u_i$ is a move slot with a target, a switch destination
or a pass. Examples are
\begin{lstlisting}
move 1 1, move 2 -1
move 1 2 mega, switch 4
\end{lstlisting}
Positive targets are opposing positions and negative targets are our own. The
opponent's simultaneous choice is \emph{not} part of a candidate. The
enumerator takes the Cartesian product of the per-position options. It
removes combinations that violate shared constraints, filters disabled or
exhausted moves, respects visible trapping, and handles fainted positions.
More than two active positions are rejected. Hidden restrictions can still
make the server refuse a candidate, so $\A(r)$ is \emph{request-visible
admissibility}, not proof of acceptance.

\begin{proposition}[Constraints enforced by joint enumeration]
\label{prop:joint}
Every emitted doubles candidate has distinct switch destinations, at most one
Mega-family activation and at most one Terastallization. In a
forced-replacement request, the number of switching positions equals
$\min(\text{forced positions},\text{available bench members})$.
\end{proposition}
\begin{proof}
Each product element is tested in turn. Its switch strings must form a set,
so duplicates are rejected. Under forced switching, its switch count must
equal the stated minimum. Elements with more than one Mega-family event or
more than one Terastallization event are rejected. An emitted element has
passed every predicate.
\end{proof}

The proposition covers these predicates only. Champions requests never offer
Terastallization, so that branch is dormant. Team validation and other hidden
mechanics are outside its scope.

\subsection{Team preview and an exact symmetry}
With $n$ party members and $k$ to bring, preview candidates are ordered
selections, $|\A_{\mathrm{preview}}|=n!/(n-k)!$, which is 360 for six choose
four. The first two positions lead; the last two are the bench. Under
\code{opening-v2} and every \code{strategic-*} profile except
\code{strategic-turn-v1}, the action encoder
canonicalizes the bench by sorting its indices before hashing. The
representation is therefore exactly invariant to swapping the two bench
positions. The actor still scores all 360 candidates, so each plan appears
twice.

\begin{proposition}[Bench symmetry halves the preview space]
\label{prop:preview}
Let $\sim$ identify two preview candidates that agree in their ordered leads
and in their bench as a set. For six choose four this gives 180 classes of
exactly two candidates each. If the features, and hence the logits, are
constant on each class, then:
\begin{enumerate}
  \item the class probabilities $\Pi(C)=\sum_{a\in C}\pi(a)$ are a softmax
    over the 180 class logits, at the same temperature;
  \item the PPO ratio of any candidate equals the ratio of its class;
  \item the candidate-level entropy is $H(\pi)=H(\Pi)+\log2$.
\end{enumerate}
\end{proposition}
\begin{proof}
There are $6\cdot5=30$ ordered lead pairs. Each leaves four members, from
which $\binom{4}{2}=6$ bench sets can be formed, so there are 180 classes.
Each bench set has exactly two orderings, which accounts for all
$180\cdot2=360$ candidates. Let $\ell_C$ be the shared logit. Then
$\pi(a)=e^{\ell_C/\tau}/\sum_{C'}2e^{\ell_{C'}/\tau}$ for $a\in C$, so
$\Pi(C)=2\pi(a)=e^{\ell_C/\tau}/\sum_{C'}e^{\ell_{C'}/\tau}$. The ratio
$\pi_\theta(a)/\pi_{\mathrm{old}}(a)=\Pi_\theta(C)/\Pi_{\mathrm{old}}(C)$
because the factor~2 cancels. Finally,
$H(\pi)=-\sum_C2\pi_C\log\pi_C=-\sum_C\Pi(C)\log(\Pi(C)/2)=H(\Pi)+\log2$.
\end{proof}

\intuition{The duplication is harmless for learning. Ratios are unchanged, and
the entropy bonus is shifted by a constant whose gradient is zero. Statements
such as ``the policy chooses among 360 previews'' should still say 180 distinct
plans. Lead order is kept on purpose: it fixes which Pok\'{e}mon occupies the
left and right positions.}

\subsection{Why joint scoring rather than a product policy}
Sampling each position independently gives
$\pi(u_1,u_2\mid x)=p(u_1\mid x)q(u_2\mid x)$. That family cannot express
coordination.

\begin{proposition}[A coordination a product policy cannot represent]
\label{prop:product}
With two options $A,B$ per position, no unmasked product distribution gives
positive mass to both $(A,A)$ and $(B,B)$ and zero mass to both $(A,B)$ and
$(B,A)$.
\end{proposition}
\begin{proof}
Positive diagonal mass forces $p(A),q(A),p(B),q(B)>0$, so the off-diagonal
products are positive too. Equivalently, every product table has rank one and
determinant zero. The required table has determinant
$\pi(A,A)\pi(B,B)>0$.
\end{proof}

The implemented alternative scores each feasible complete candidate directly.
Masking a product policy with joint constraints, or factorizing it
autoregressively, would also introduce dependence. Feature collisions and
finite capacity still limit which joint distributions the network can express.

\subsection{Two enumerators}
The search engine enumerates candidates separately, inside the simulator, in
\code{sideChoices} of \code{search/engine.cjs}. Its pruning differs from
$\A(r)$:
\begin{itemize}
  \item Damaging moves aimed at the partner are removed.
  \item Targets in empty foe slots are removed while the other foe is alive.
  \item Every move of a Mega-capable Pok\'{e}mon carries the \code{mega} flag,
    so delaying Mega Evolution is never considered.
\end{itemize}
Our own live choices still start from $\A(r)$ and are pruned in Python. The
engine enumeration is used for the opponent and for our forced-replacement
look-ahead. Appendix~\ref{app:defects} records a defect in it: when both
actives on one side can Mega Evolve, every move-and-move pair is discarded.

\source{\code{legal_choices} in \code{battle_state.py}; bench canonicalization
in \code{Features.action} (\code{ml/features.py}); \code{_prune_ours} in
\code{search/agent.py}; \code{sideChoices} in \code{search/engine.cjs}.}
